\subsection{Selected explicit equations}
\begin{center}\small
\begin{tabular}{lrrrrr}\toprule
Domain and profile & Variables & Degree & Rows & Terms & Height\\\midrule
$\N$, dimension &206&4&158&8210&130686376\\
$\N$, balanced &207&4&159&5278&69238087\\
$\N$, sparse &208&4&160&5198&65866712\\\midrule
$\Z$, dimension &239&4&175&12658&6144\\
$\Z$, balanced &240&4&176&5232&4096\\
$\Z$, sparse &240&4&176&4927&7596\\\bottomrule
\end{tabular}
\end{center}
The natural 206-variable equation improves the matched preceding
219-variable/171-row/8267-term/height-130686428 equation in all four
varying metrics. The sparse natural point similarly improves the
preceding 221-variable/173-row/5253-term/height-65867916 equation.
The middle point provides a separate dimension/support/height tradeoff.

The new 239-variable integer quartic improves the incoming shared-error
240-variable/177-row/21482-term/height-6144 point in dimension, rows and
support at the same height. The sparse integer point improves the
incoming factorial-free 242-variable/178-row/5018-term/height-7968
comparison in every varying metric. The balanced point trades support
for a smaller coefficient height; these comparisons are not global
variable-count records.

\begin{center}
\begin{tabular}{lrrr}\toprule
Domain and circuit & Variables & Exact degree & Binary operations\\\midrule
Strictly positive, inverse correction &10&304&85\\
Unrestricted, completed-square carrier &13&152&119\\\bottomrule
\end{tabular}
\end{center}
The positive circuit improves the incoming degree-304/87-operation
Lucas-correction form at unchanged degree and domain. Its
84-operation/degree-410 square-weight companion is a dominated control.
The 13-integer circuit improves the matched degree-152/126-operation
incoming form. It remains a tradeoff against our earlier degree-156/
117-operation point. Its coordinate shear is an exact integer bijection,
not an identity in the original unshifted coordinate names.

Both numerical specializations of the positive circuit are supplied:
the original radix67 tuple and the complete previously reviewed
incoming complemented-label radix47 tuple. The latter has 207214 raw
coefficient bytes instead of 295388; its encoding is credited prior
work, not a new encoding theorem in this report. Counts and degrees
are unchanged, while the specialized integer polynomials differ.
